\begin{enumerate}
  \item L'acide acétylsalicylique ou aspirine peut être synthétisé au laboratoire à partir de la réaction entre l'acide salicylique et l'anhydride éthanoïque en utilisant le chauffage à reflux selon l'équation de la réaction suivante modélisant cette transformation :\\
\includegraphics[width=\textwidth]{2025_10_31_838bdd6d133f8b78fc82g-1}\\
1.1. Donner le nom du groupement fonctionnel délimité par un trait pointillé fermé dans la forme topologique de chacune des molécules d'acide salicylique et d'acide acétylsalicylique.\\
1.2. Citer les deux caractéristiques de cette transformation.\\
1.3. Choisir, parmi les montages expérimentaux (1), (2) et (3) le montage utilisé pour réaliser cette synthèse.\\
\includegraphics[width=\textwidth]{2025_10_31_838bdd6d133f8b78fc82g-1(1)}\\
1.4. Quel est l'intérêt du chauffage à reflux ?\\
1.5. On introduit dans une fiole jaugée, $n_{1}=0,10 \mathrm{~mol}$ d'acide salicylique et $n_{2}=0,26 \mathrm{~mol}$ d'anhydride éthanoïque et quelques gouttes d'acide sulfurique concentré. Après chauffage à reflux, et les opérations de traitement et de purification, on obtient des cristaux d'aspirine de masse $m_{\text {exp }}=15,3 \mathrm{~g}$.\\
Calculer le rendement de cette synthèse sachant que le réactif limitant est l'acide salicylique.\\
On donne : Masse molaire de l'acide acétylsalicylique : $M=180 \mathrm{~g} \cdot \mathrm{~mol}^{-1}$
  \item On prépare une solution aqueuse ( $S$ ) d'acide acétylsalicylique de concentration molaire\\
$C=5,55.10^{-3} \mathrm{~mol} . \mathrm{L}^{-1}$ et de volume $V=500 \mathrm{~mL}$. Après mesure de la conductivité de la solution $(S)$, on détermine la valeur de $x_{\mathrm{f}}$ avancement de la réaction à l'état final du système chimique, et on trouve: $x_{\mathrm{f}}=5,70 \cdot 10^{-4} \mathrm{~mol}$.\\
Pour simplifier, on désigne la molécule de l'acide acétylsalicylique par $A H$, et sa base conjuguée par $A^{-}$.\\
2.1. Ecrire l'équation de la réaction de l'acide acétylsalicylique avec l'eau.\\
2.2. Montrer que la réaction de l'acide acétylsalicylique avec l'eau est non totale.\\
2.3. Déterminer la valeur de la constante d'acidité $K_{\mathrm{A}}$ du couple $A H_{(\mathrm{aq})} / A_{(\mathrm{aq})}^{-}$.
\end{enumerate}

